\section{Python: essentials}
Python 3.10+; standard library only. Run: \texttt{python3 sol.py < 1.in}.

\subsection{Input and output}
\begin{lstlisting}[label=py:io]
n = int(input())
a = list(map(int, input().split()))
print(sum(a))
print(*a)                         # space-separated
\end{lstlisting}
\hint{Built-in \texttt{input()} returns a string without its final newline.
Keep it as the default. Optional faster line input still returns strings:}
\begin{lstlisting}
# import sys
# input = sys.stdin.readline  # optional; keeps newline
\end{lstlisting}
\hint{Optional whole-file integer input (not interactive; uses more memory):}
\begin{lstlisting}[label=py:tokens]
import sys
it = iter(map(int, sys.stdin.buffer.read().split()))
n = next(it)
a = [next(it) for _ in range(n)]
\end{lstlisting}
\hint{\texttt{buffer} gives bytes; use \texttt{token.decode()} for text.
Interactive: \texttt{print("?", x, flush=True)}.}

\subsection{Control flow and functions}\label{py:syntax}
\begin{lstlisting}
def bounds(a):
    if not a:
        return None
    return min(a), max(a)        # tuple return

for i, x in enumerate([3, 1, 4]):
    if x < 0:
        continue
    if x == 4:
        break
else:
    print("no break")

x = 20
while x > 0:
    x //= 2
\end{lstlisting}
\hint{Function assignment creates a local variable. Use \texttt{nonlocal}
for a name in an enclosing function, \texttt{global} for a module name.
Mutating a passed list changes that list; rebinding the parameter does not.
Use \texttt{def f(a=None)} and create the list inside; \texttt{a=[]} as a
function default reuses the same list across calls.}

\subsection{Input formats and output batches}\label{py:formats}
\begin{lstlisting}
t = int(input())                 # only if t is specified
for _ in range(t):
    n = int(input())
    a = list(map(int, input().split()))

# Alternative: read lines until EOF (no test count).
import sys
for line in sys.stdin:
    if not line.strip():
        continue
    a = list(map(int, line.split()))

out = [str(x*x) for x in range(5)]
print("\n".join(out))
\end{lstlisting}
\hint{These are separate input formats: follow the statement.
Do not assume one integer per line. For a fixed number of tokens spread
across lines, use the optional whole-file reader above.
Interactive: read one reply per question and flush each question.}

\subsection{Arithmetic and comparisons}
\begin{lstlisting}[label=py:arithmetic]
7 / 3                  # float: 2.333...
-7 // 3, -7 % 3        # (-3, 2): division floors
a, b = -7, 3           # b > 0
ceil_div = -(-a // b)  # -2, works for either sign of a
2 ** 10                # 1024; ^ is XOR, not power
10 ** 18               # exact int; 1e18 is a float
x, y = 3, 8; x, y = y, x  # swap; unpack a tuple
0 <= x < 10            # chained comparison
x == y                 # value equality; "is": identity
\end{lstlisting}

\subsection{Lists, slices and grids}
\begin{lstlisting}[label=py:lists]
a = [10, 20, 30, 40]
a.append(50); a.extend([60, 70])
last = a.pop()                    # removes last
a[-1], a[1:3], a[::-1]           # 60, [20,30], reversed
b = a.copy()                     # shallow copy
a[1:3] = [9]                     # replace a slice
n, m = 3, 4
grid = [[0] * m for _ in range(n)]
squares = [x*x for x in range(6) if x % 2 == 0]
\end{lstlisting}
\hint{\texttt{a=b} aliases; \texttt{[[0]*m]*n} repeats the same row.
Slices copy and exclude the stop index.}

\subsection{Loops and built-ins}
\begin{lstlisting}[label=py:builtins]
a = [3, 1, 4]
list(range(2, 8, 2))              # [2, 4, 6]
list(enumerate(a))                # [(0,3),(1,1),(2,4)]
list(zip(a, [10, 20]))            # stops at shorter input
len(a), sum(a), min(a), max(a)    # 3, 8, 1, 4
any(x < 0 for x in a)             # False
all(x > 0 for x in a)             # True
list(reversed(a))                # [4, 1, 3]
\end{lstlisting}
\hint{\texttt{map/zip/reversed} are one-pass iterators;
\texttt{list(...)} stores their results.}

\section{Python: exact numbers}

\subsection{Huge integers and modular arithmetic}
\begin{lstlisting}[label=py:integers]
a = 10 ** 100                     # arbitrary-size int
q, r = divmod(a, 97)              # a == q*97 + r
int("ff", 16), int("101", 2)      # 255, 5
bin(13), hex(255)                 # '0b1101', '0xff'
pow(2, 1000, 97)                  # 2**1000 mod 97
pow(3, -1, 10)                    # 7: inverse, gcd=1
abs(-123), len(str(a))            # 123, 101 digits
\end{lstlisting}
\hint{\texttt{pow(a,e,m)} avoids building \texttt{a**e}.
Inverses need coprime \texttt{a,m}; \texttt{m} need not be prime.}

\hint{For extremely large decimal input/output, CPython may limit
\texttt{int}/\texttt{str} conversion (often 4300 digits). For a trusted contest
problem requiring more digits, optionally disable that conversion limit:}
\begin{lstlisting}
import sys
if hasattr(sys, "set_int_max_str_digits"):
    sys.set_int_max_str_digits(0)
\end{lstlisting}

\subsection{Integer math and counting}
\begin{lstlisting}[label=py:math]
from math import gcd, lcm, isqrt, comb, perm, factorial
gcd(12, 18, 30), lcm(4, 6, 10)    # 6, 60
isqrt(10 ** 100 + 1)              # 10**50, exact floor
comb(100, 50)                     # exact n choose k
perm(5, 2), factorial(5)          # 20, 120
from math import inf, pi, hypot, atan2, floor, ceil
hypot(3, 4), atan2(1, 0)          # 5.0, pi/2 radians
\end{lstlisting}
\hint{\texttt{comb/factorial} give full integers, not values modulo
\texttt{MOD}. Use \texttt{isqrt}, not float \texttt{sqrt}, for exact roots.}

\subsection{Exact fractions}
\begin{lstlisting}[label=py:fractions]
from fractions import Fraction as F
x = F(1, 3) + F(1, 6)            # Fraction(1, 2)
x.numerator, x.denominator       # (1, 2)
F("0.125")                       # 1/8, exact decimal
F(2, 3) < F(3, 4)                # exact comparison
F("3.14159").limit_denominator(100)  # 311/99
\end{lstlisting}
\hint{Fractions reduce automatically. Use integers or strings:
\texttt{F(0.1)} preserves a binary float, not exactly \texttt{1/10}.}

\subsection{High-precision decimal arithmetic}
\begin{lstlisting}[label=py:decimal]
from decimal import Decimal as D, localcontext
with localcontext() as ctx:
    ctx.prec = 60                 # significant digits
    x = D(1) / D(7)              # rounded to precision
    root = D(2).sqrt()
    y = D("0.1") + D("0.2")      # exactly 0.3
    z = D("1.235").quantize(D("0.01"))  # 1.24
\end{lstlisting}
\hint{Use strings for decimal input. Context controls precision and
rounding. For exact \texttt{1/3}, use \texttt{Fraction}.}

\subsection{Integer bitmasks}
\begin{lstlisting}[label=py:bits]
mask, i, n = 0b10110, 1, 6
mask.bit_count(), mask.bit_length()  # 3, 5
bool(mask & (1 << i))             # test bit i
mask |= 1 << 5                    # set bit 5
mask &= ~(1 << 1)                 # clear bit 1
mask ^= 1 << 2                    # toggle bit 2
mask << 3, mask >> 1              # unlimited-width shifts
(~mask) & ((1 << n) - 1)          # complement in n bits
\end{lstlisting}

\section{Python: containers}

\subsection{Dictionaries and composite keys}
\begin{lstlisting}[label=py:dict]
d = {"a": 2}
d["b"] = 3
d.get("missing", 0)               # 0; does not insert
"a" in d                         # membership tests keys
list(d.items())                  # [('a',2),('b',3)]
list(d.keys()), list(d.values())
d.pop("missing", None)            # default avoids KeyError
positions = {(2, 3): 7}          # tuple key, built-in hash
frozen = frozenset([1, 2, 2])     # hashable set: {1,2}
\end{lstlisting}
\hint{Dict iteration preserves insertion order. Keys must be hashable:
use tuples/frozensets, not lists/sets. Missing \texttt{d[k]} raises
\texttt{KeyError}.}

\subsection{Sets}
\begin{lstlisting}[label=py:sets]
s, t = {1, 2, 3}, {3, 4}
s | t, s & t                     # union, intersection
s - t, s ^ t                     # difference, symmetric diff
s <= t, s.isdisjoint(t)           # subset, no overlap
s.add(5); s.discard(99)           # discard tolerates missing
empty = set()                    # {} is an empty dict
unique = set([3, 1, 3])
\end{lstlisting}
\hint{Sets have no sorted-order guarantee.
Set/dict membership is average $O(1)$; list membership is $O(n)$.}

\subsection{Frequency tables and multisets}
\begin{lstlisting}[label=py:counter]
from collections import Counter
c = Counter("banana")            # a:3, n:2, b:1
c["x"]                          # 0; does not insert
c.most_common(2)                 # [('a',3), ('n',2)]
list(c.elements())               # repeat by multiplicity
c + Counter("an"), c - Counter("an")
c & Counter("band"), c | Counter("band")
# &: minimum counts; |: maximum; +/-: add/subtract
\end{lstlisting}
\hint{Counter arithmetic drops nonpositive counts.
\texttt{c.subtract(other)} mutates and keeps zero/negative counts.}

\subsection{Grouping with missing defaults}
\begin{lstlisting}[label=py:defaultdict]
from collections import defaultdict
groups = defaultdict(list)
groups["a"].append(4)             # creates [] if missing
counts = defaultdict(int)
counts["a"] += 1                 # creates 0 if missing
sets = defaultdict(set)
sets["a"].add(4)
\end{lstlisting}

\subsection{Deque: queue and both ends}
\begin{lstlisting}[label=py:deque]
from collections import deque
q = deque([1, 2])                 # O(1) ends; list.pop(0): O(n)
q.append(3); q.appendleft(0)
q.popleft(), q.pop()              # 0, 3
q.rotate(1)                      # [2,1]: right by 1
q[0], q[-1]                      # 2, 1
\end{lstlisting}



\section{Python: strings and ordering}

\subsection{Strings}
\begin{lstlisting}[label=py:strings]
s = "  one  two\n"
s.strip(), s.split()              # 'one  two', ['one','two']
"a,,b,".split(",")                # ['a','','b','']
"-".join(["a", "b", "c"])         # 'a-b-c'
s.replace("one", "1")             # replace all occurrences
"banana".find("na")               # 2; -1 if absent
"banana".count("an")              # 2, non-overlapping
"abc".startswith("ab"), "abc".endswith("c")
"abc"[::-1], "AbC".lower()        # 'cba', 'abc'
ord("a"), chr(97)                 # 97, 'a'
\end{lstlisting}
\hint{Strings are immutable. Build many output fragments in a list,
then join; repeated concatenation can be quadratic.
\texttt{split()} collapses whitespace; \texttt{split(",")} keeps
empty fields.}

\subsection{Formatting and conversions}
\begin{lstlisting}[label=py:format]
x = 1 / 3
f"{x:.10f}"                      # '0.3333333333'
f"{7:03d}", f"{255:x}"           # '007', 'ff'
f"{13:08b}"                      # '00001101'
"  abc  ".strip()                # 'abc'
int("-42"), float("1.25")        # -42, 1.25
str(123), list("abc")             # '123', ['a','b','c']
\end{lstlisting}

\subsection{Sorting and key functions}
\begin{lstlisting}[label=py:sorting]
a = [(2, 5), (1, 7), (2, 3)]
b = sorted(a)                    # lexicographic tuples
a.sort(key=lambda p: (p[0], -p[1]))
sorted(["bbb", "a", "cc"], key=len)
sorted([3, 1, 2], reverse=True)   # [3,2,1]
max(a, key=lambda p: p[1])        # item with largest second
min([], default=0)               # 0 instead of ValueError
\end{lstlisting}
\hint{Sorts are stable; equal keys keep their original order.
\texttt{a.sort()} mutates and returns \texttt{None};
\texttt{sorted(a)} makes a new list.}

\subsection{Binary search in a sorted list}
\begin{lstlisting}[label=py:bisect]
from bisect import bisect_left as bl, bisect_right as br
from bisect import insort
a, x = [1, 2, 2, 5], 2
bl(a, x), br(a, x)                # 1, 3
count = br(a, x) - bl(a, x)       # 2 occurrences
insort(a, 3)                     # [1,2,2,3,5]
\end{lstlisting}
\hint{Search: $O(\log n)$. Insertion: $O(n)$.
This is not a balanced ordered set.}

\subsection{Extracting structured text}
\begin{lstlisting}[label=py:regex]
import re
re.findall(r"-?\d+", "x=-12, y=34")  # ['-12','34']
re.fullmatch(r"[A-Z]+\d+", "AB12") is not None
re.split(r"[,;]+", "a,b;;c")      # ['a','b','c']
\end{lstlisting}

\section{Python: iteration and helpers}

\subsection{Combinatorial iterators}
\begin{lstlisting}[label=py:itertools]
from itertools import (combinations, permutations, product,
    combinations_with_replacement, accumulate, chain, pairwise)
list(combinations("ABC", 2))     # AB, AC, BC (as tuples)
list(permutations("ABC", 2))     # AB, AC, BA, BC, CA, CB
list(product([0, 1], repeat=3))  # 000, 001, ..., 111
list(product("AB", [1, 2]))      # A1, A2, B1, B2
list(combinations_with_replacement("AB", 2))  # AA, AB, BB
list(accumulate([2, 3, 4], initial=0))  # [0,2,5,9]
list(chain.from_iterable([[1, 2], [3]]))  # [1,2,3]
list(pairwise([4, 7, 9]))        # [(4,7),(7,9)]
\end{lstlisting}
\hint{These return lazy iterators; \texttt{list(...)} above shows
the results. Iterate directly to avoid storing all outputs.
Repeated input values can give duplicate permutations.}

\subsection{Automatic memoization}
\begin{lstlisting}[label=py:cache]
from functools import cache

@cache
def f(i, mask):
    ...                          # translate the recurrence

f.cache_clear()                  # between different test cases
# Older Python: @lru_cache(None), from functools import lru_cache
\end{lstlisting}
\hint{Arguments must be hashable (e.g.\ ints, tuples, frozensets).
Cache keys include arguments, not changing global variables.
Clear between cases or make a new function. Recursion depth still
applies; iterative traversal is safer for deep graphs.}

\subsection{Min-heaps and sorted merging}
\begin{lstlisting}[label=py:heap]
from heapq import (heapify, heappush, heappop, heappushpop,
                  heapreplace, nsmallest, nlargest, merge)
h = [5, 2, 8]
heapify(h)                       # O(n); minimum at h[0]
heappush(h, 3)
heappop(h)                       # 2; O(log n)
heappushpop(h, 4)                # push, then pop minimum
heapreplace(h, 6)                # pop minimum, then push
nlargest(2, [4, 1, 7]), nsmallest(2, [4, 1, 7])
list(merge([1, 4], [2, 3]))      # [1,2,3,4]; inputs sorted
\end{lstlisting}
\hint{Min-heap: use \texttt{(-priority, item)} for a max-heap.
Tuples compare lexicographically; add a unique counter before
non-comparable payloads. \texttt{heapreplace} needs a nonempty
heap. Merging is lazy.}

\subsection{Comprehensions and generators}
\begin{lstlisting}[label=py:comprehensions]
a = [3, 1, 3]
index = {x: i for i, x in enumerate(a)}  # last index of x
pairs = {(x, y) for x in range(2) for y in range(3)}
total = sum(x*x for x in a)      # generator: no extra list
first, *middle, last = [1, 2, 3, 4]
\end{lstlisting}
\hint{Use brackets for a list, braces for a set/dict,
\texttt{(expr for ...)} for a generator.}

\section{Python: dates and reminders}

\subsection{Calendar arithmetic}
\begin{lstlisting}[label=py:dates]
from datetime import date, timedelta
a, b = date(2024, 2, 28), date(2024, 3, 1)
(b - a).days                     # 2 (leap year)
a + timedelta(days=2)            # 2024-03-01
a.weekday()                      # 2: Monday=0, Sunday=6
a.isoformat()                    # '2024-02-28'
date.fromisoformat("2024-03-01")  # date(2024,3,1)
from calendar import isleap, monthrange
isleap(2024), monthrange(2024, 2)[1]  # True, 29 days
\end{lstlisting}
\hint{\texttt{date} supports years 1--9999.
\texttt{monthrange(y,m)} returns (weekday of day 1, days in month).}

\subsection{Topological order}
\begin{lstlisting}[label=py:topological]
from graphlib import TopologicalSorter, CycleError
pred = {"A": set(), "B": {"A"}, "C": {"A", "B"}}
order = list(TopologicalSorter(pred).static_order())
# ['A','B','C']; CycleError if the graph has a cycle
\end{lstlisting}
\hint{Input lists \textbf{predecessors}, not outgoing neighbours.
With several valid orders, the returned order need not be
lexicographically smallest.}

\subsection{JSON input}
\begin{lstlisting}[label=py:json]
import json
obj = json.loads('{"a": [1, 2], "ok": true}')
obj["a"][0], obj["ok"]           # 1, True
json.dumps(obj)                  # encode back to JSON
\end{lstlisting}

\subsection{C++ to Python: avoid surprises}
\begin{itemize}
\item Indentation defines blocks; use 4 spaces.
\item \texttt{and/or/not}, \texttt{True/False/None}; no \texttt{++}.
\item \texttt{append/sort/reverse} mutate and return \texttt{None}.
\item \texttt{==} compares values; \texttt{is} tests identity.
      Use \texttt{x is None} for a missing value.
\item \texttt{(x,)} is a one-element tuple; \texttt{(x)} is just \texttt{x}.
\item Iterating a dict gives keys. Do not change a dict/set's size
      while iterating it.
\item Integers do not overflow; floats still round.
\end{itemize}

\subsection{Complexities to remember}\label{py:complexities}
\small
\begin{tabularx}{\linewidth}{@{}Xr@{}}
\toprule
Operation & Cost\\
\midrule
List index / last append or pop & $O(1)$ amort.\\
List search / front or middle insert & $O(n)$\\
Dict/set lookup, add, remove & $O(1)$ avg.\\
Deque append/pop at either end & $O(1)$\\
Heap push/pop & $O(\log n)$\\
Sort & $O(n\log n)$\\
List/string slice of length $k$ & $O(k)$\\
\bottomrule
\end{tabularx}
\normalsize
\hint{Bounds above treat numbers/keys as fixed-size.
Big integer arithmetic and hashing long keys depend on their size.}
